% discussion.tex
Pattern multiplicity in the modern Hopfield energy provides a training-free way to condition
protein sequence generation on user-designated subsets. Multiplicity weights add a logit bias
to the Boltzmann distribution. The corresponding continuous-time equilibrium target is an
exact Gaussian mixture. Finite-step ULA retains discretization and mixing error. Its decoded
marker fractions and amino acid compositions were similar to independent exact-equilibrium
draws. Exact sampling also reproduced the large gap between
$f_{\mathrm{eff}}=0.996$ and decoded P1~K/R recovery at $\rho=500$, indicating that
finite-step ULA error was not the main source of that gap.
A single scalar $\rho$ continuously interpolates between unconditioned generation
and hard subset curation. The required $\rho$ for a target designated fraction follows
directly from the family composition (Eq.~\ref{eq:rho-inverse}). In the Kunitz domain,
curated memory retained the selected K/R-at-P1 marker with as few as three designated
inputs. Full-family multiplicity weighting increased marker fidelity from 32\% to
\KunitzFobsFiveHundred{} $\pm$ \KunitzFobsFiveHundredSD{} at $\rho = 500$.
Results across five Pfam families showed an association between PCA-space geometry and
conditioning effectiveness.

We observed a calibration gap between latent designated mass and decoded marker frequency.
The mean attention differed from $f_{\mathrm{eff}}$ by at most \AttnMaxDevPts{} percentage
points across the canonical sweeps. The remaining gap arises as latent states pass through
finite-temperature noise, unit-normalized PCA coordinates, affine reconstruction, and argmax
decoding (Eq.~\ref{eq:decoder-pushforward}). When the designated and background subsets
overlap in PCA space, the energy bias may not change the decoded residue identity. Across five
Pfam families, the exploratory linear fit had
$R^{2}=\RelationRsq$, but it was sensitive to individual families: leave-one-family-out slopes
ranged from \RelationLooSlopeMin{} to \RelationLooSlopeMax{} and $R^{2}$ from
\RelationLooRsqMin{} to \RelationLooRsqMax{}. The separation index should therefore be treated
as a descriptive diagnostic rather than an \emph{a priori} predictor or decision threshold.
Broader prospective validation is required before it can guide the choice between multiplicity
weighting and hard curation.

Hard curation retained the marker with as few as three designated input sequences. Small
characterized sets can arise from screens, functional selections, or literature curation.
SA can use such a set to generate additional sequences that share its sequence features.
In the Kunitz experiment, selection and evaluation both used K/R at P1. Complete recovery
therefore shows marker retention rather than functional validation. Generated candidates would
still require experimental screening. The method does not assign a biological meaning to the
designation. The scaling study shows that diversity increased with the number of designated
inputs and changed little beyond $K_{\mathrm{des}}=20$.

The $\omega$-conotoxin experiment provides a second example of conditioning from a small
designated subset. It also compares full-family and curated-subset seeding. Full-family seeding is appropriate when
the goal is to generate novel family members with broad structural diversity; curated seeding is
necessary when the goal is to preserve a specific sequence marker. In the $\omega$-conotoxin case,
the O-superfamily is structurally coherent (all members share the C--C--CC--C--C disulfide
framework) but functionally heterogeneous (members block N-type, P/Q-type, or sodium channels
with varying selectivity). Full-family seeding dilutes the Tyr13-associated marker
because its PCA encoding reflects the full sequence diversity. Curated seeding restricts the
PCA basis to variation among the 23 designated accessions. It recovered Tyr13 at 98\%
frequency. The accession designation and Tyr13 marker agree for only
86.5\% of stored sequences, so these outcomes are related but not interchangeable. This illustrates
the separation-index interpretation: even
when $S$ is moderate, hard curation can achieve high marker retention by
restricting the memory to the designated subset. The cost (lower sequence diversity
and reduced novelty, 0.46 vs.\ 0.60) reflects the narrower PCA basis, consistent with the
Kunitz scaling study.

The AlphaFold2-multimer analysis produced low-confidence predictions for generated and
control sequences: iPTM was approximately 0.10 and interface pLDDT approximately 37--41.
Permutation tests detected no group differences, but nonsignificance at sample sizes 10/10/5
does not establish equivalence or preserved binding geometry. These calculations therefore
provide no positive evidence about Cav2.2 interaction; experimental complex structures and
binding assays are needed.

Several limitations should be acknowledged. The cross-family regression spans only five Pfam
families, with $\omega$-conotoxin displayed externally rather than included in the fit. Its
leave-one-family-out sensitivity and nonmonotonic points preclude a validated predictive rule;
additional prospectively selected families are needed to characterize the $S$--$\Delta$
relationship. Our Pfam splits are defined by data-selected single marker positions
(P1 in Kunitz, Trp in SH3, a specificity-loop residue in WW, Gln at position~50 in
Homeobox, and H/N in the Forkhead recognition helix). Multi-position phenotypes may show
different calibration behavior because they depend on coordinated residue identities.
We have not experimentally
validated the binding activity of generated sequences; marker fidelity is not evidence of
functional binding. For the $\omega$-conotoxin family, we compared generated residue frequencies
with positions reported by alanine-scanning and site-directed mutagenesis studies
~\cite{kimTyr13Essential1995,satoBasicResidues1993,lewisConotoxinSAR2012}. Designated-seeded
generation retained the cysteine scaffold and Tyr13 signal and enriched several reported
positions (Table~\ref{tab:sar-agreement}), but this does not establish their effects in the new
sequence backgrounds. For the Kunitz domain, deep mutational scanning
data for the BPTI--trypsin interaction are now
available~\cite{heyneBPTIDMS2021} and could provide a quantitative binding-landscape comparison
in future work.

Looking forward, the multiplicity framework opens several directions. The binary split
(designated/background) generalizes naturally to multi-class conditioning by assigning different
multiplicities to different functional categories. The PCA bottleneck that limits conditioning
effectiveness might be reduced by alternative sequence encodings. Learned representations from
protein language models, for instance, may better separate functional subsets than linear PCA.
The $\beta$ lever suggests that adaptive temperature schedules, increasing $\beta$ during sampling
to sharpen the multiplicity bias, could improve marker transfer without permanently sacrificing
diversity. These extensions remain within the analytic, training-free framework that makes
stochastic attention practical for the scarce-data regime.
